% =============================================
% File: chapters/conclusion.tex
% KẾT LUẬN VÀ KIẾN NGHỊ
% =============================================

\chapter*{KẾT LUẬN VÀ KIẾN NGHỊ}
\addcontentsline{toc}{chapter}{Kết luận và kiến nghị}
\label{chap:ketluan}


% --------------------------------------------------
\section*{1. Kết luận}
% --------------------------------------------------

Đề tài \textit{``Nghiên cứu và ứng dụng bộ điều khiển PID kết hợp mạng nơron
nhân tạo (ANN) điều khiển Robot Bipedal Wheel sử dụng cơ cấu 5-bar parallel''}
đã được thực hiện theo hai hướng song song: \textbf{điều khiển cổ điển kết hợp
học máy} và \textbf{học tăng cường trong mô phỏng vật lý}. Các kết quả đạt được:

\begin{enumerate}[leftmargin=1.5cm, label=\arabic*.]
    \item \textbf{Nhận dạng hàm truyền động cơ NF5475.}
    Thực nghiệm trên sáu mẫu động cơ, sử dụng công cụ \texttt{tfest} (MATLAB)
    với cấu trúc bậc 4/3. Hàm truyền đại diện $G_1(s)$ đạt độ chính xác
    95{,}36\%, hệ số khuếch đại DC trung bình $\bar{K}_{DC} = 14{,}23$
    với độ phân tán $\pm 2\%$ qua 6 mẫu — xác nhận tính đồng đều của lô
    động cơ và phản ánh các hiệu ứng hao mòn chổi than theo tuổi thọ sử dụng.

    \item \textbf{Tối ưu bộ điều khiển cascade PID bằng giải thuật di truyền (GA).}
    GA với 80 cá thể, 80 thế hệ, 20 lần chạy độc lập và hàm phạt ràng buộc
    cho phép tìm bộ tham số $\{K_{pv}^*, K_{iv}^*\}$ tại 9 mức setpoint
    (250--2250~RPM). Kết quả: overshoot $\leq 5\%$, thời gian xác lập
    $\leq 0{,}65\,$s, sai số xác lập $< 4\,$RPM — vượt trội so với Random
    và MATLAB PID Tuner.

    \item \textbf{Xây dựng mạng ANN dự đoán tham số PI thích nghi.}
    Mạng $2 \to 8 \to 16 \to 2$ (ReLU, Dropout, L2-regularization)
    được huấn luyện thuần túy bằng Python (không dùng framework),
    hội tụ sau $\approx 400$ epochs trên 9 cặp dữ liệu GA.
    Mạng học được quan hệ đơn điệu setpoint $\to$ tham số PI và
    có khả năng nội suy cho setpoint chưa thấy trong huấn luyện.

    \item \textbf{Triển khai bộ điều khiển PI trên vi điều khiển ESP32 WROOM-32.}
    Vòng PI rời rạc (Tustin, $\Delta t = 10\,$ms) với anti-windup hoạt động
    ổn định ở toàn dải setpoint thử nghiệm, đáp ứng vận tốc đạt trong
    $\pm 5\%$ sau $t_{rise} < 0{,}5\,$s.

    \item \textbf{Phân tích động học cơ cấu 5-bar parallel của robot Legged~V3.}
    Xác định động học thuận, động học nghịch và ma trận Jacobian với thông số
    thực tế ($A_1A_2 = 105\,$mm, $L_{act} = 100\,$mm, $L_{pas} = 180\,$mm),
    làm cơ sở cho điều khiển quỹ đạo điểm cuối chân.

    \item \textbf{Huấn luyện chính sách điều khiển toàn thân bằng PPO trong Isaac Lab.}
    Môi trường mô phỏng song song (4\,096 env, GPU RTX~3070) được xây dựng
    với cơ cấu 5-bar, IMU, encoder bánh xe và curriculum 3~pha dựa trên EMA
    hiệu suất. Chính sách Actor--Critic bất đối xứng ([256,256,256] ELU)
    học phối hợp 4~DOF (2~khớp hông + 2~bánh xe) để bám lệnh vận tốc
    và duy trì chiều cao thân. Trọng số Actor được xuất sang ONNX
    chuẩn bị triển khai thực tế.
\end{enumerate}

Những kết quả trên chứng minh tính khả thi của việc kết hợp hai hướng nghiên
cứu: điều khiển cổ điển tối ưu bằng GA với ANN thích nghi (phù hợp phần cứng
nhúng chi phí thấp) và học tăng cường trong mô phỏng vật lý quy mô lớn
(phù hợp robot có động học phức tạp như cơ cấu 5-bar). Cả hai cung cấp
góc nhìn bổ trợ và tạo nền tảng cho robot hai chân có bánh xe hoạt động tự
chủ trong thực tế.


% --------------------------------------------------
\section*{2. Hạn chế}
% --------------------------------------------------

\begin{enumerate}[leftmargin=1.5cm, label=\arabic*.]
    \item \textbf{Chưa kiểm chứng ANN trực tuyến trên phần cứng.}
    Mạng ANN hiện chỉ được đánh giá offline; chưa nhúng lên ESP32 để đo
    thời gian suy luận và độ chính xác dự đoán trong điều kiện thời gian thực.

    \item \textbf{Thực nghiệm động cơ ở tải cố định.}
    Thực nghiệm thu thập dữ liệu và kiểm chứng PI được thực hiện với tải
    trục cố định; chưa khảo sát hiệu năng khi tải thay đổi đột ngột
    hoặc có nhiễu ngoài (rung, va chạm).

    \item \textbf{Khoảng cách mô phỏng--thực tế (sim-to-real gap) chưa được
    kiểm chứng.}
    Chính sách PPO được huấn luyện trong Isaac Lab với vật lý PhysX lý tưởng;
    chưa thực nghiệm triển khai ONNX lên phần cứng robot thực để đánh giá
    ảnh hưởng của ma sát, độ trễ cảm biến và sai số mô hình.

    \item \textbf{Mô hình ANN chỉ bao phủ một động cơ đơn.}
    Chưa mở rộng sang điều khiển đồng thời hai động cơ trong cơ cấu 5-bar
    với ràng buộc động học chuỗi kín.

    \item \textbf{Chỉ có termination ``time\_out'' trong môi trường RL.}
    Các điều kiện kết thúc sớm (góc nghiêng quá lớn, tiếp xúc bất hợp lệ)
    đang bị tắt để tạo điều kiện khám phá; điều này có thể dẫn đến chính
    sách chưa đủ mạnh khi gặp nhiễu lớn trong thực tế.
\end{enumerate}


% --------------------------------------------------
\section*{3. Kiến nghị và hướng phát triển}
% --------------------------------------------------

\begin{enumerate}[leftmargin=1.5cm, label=\arabic*.]
    \item \textbf{Nhúng ANN lên ESP32 và kiểm chứng real-time.}
    Lượng tử hóa hoặc chuyển đổi mô hình sang định dạng nhẹ (fixed-point),
    nhúng lên ESP32 và đo thời gian suy luận, so sánh đầu ra với bảng tra GA.

    \item \textbf{Triển khai ONNX Actor lên phần cứng robot thực.}
    Chạy suy luận ONNX trên máy tính nhúng (Raspberry Pi / Jetson Nano) kết
    nối với ESP32 qua UART, đánh giá sim-to-real gap và áp dụng kỹ thuật
    domain randomization hoặc system identification thực tế nếu cần.

    \item \textbf{Bật lại điều kiện kết thúc sớm trong RL.}
    Kích hoạt \texttt{bad\_orientation} và \texttt{illegal\_contact} sau khi
    chính sách đã ổn định ở Pha~2, kết hợp với domain randomization
    (ma sát, khối lượng, độ trễ) để tăng tính bền vững.

    \item \textbf{Mở rộng bộ điều khiển PI--ANN sang hai động cơ.}
    Tích hợp IK và Jacobian 5-bar để điều khiển quỹ đạo điểm cuối chân,
    bổ sung điều khiển hai động cơ trong mỗi cụm chân.

    \item \textbf{Tích hợp cảm biến thực tế vào vòng phản hồi.}
    Gắn IMU thực (ví dụ MPU-6050) lên thân robot, xử lý tín hiệu và đưa
    vào vòng điều khiển thay thế mô phỏng, chuẩn bị cho điều khiển cân bằng
    động toàn thân.

    \item \textbf{Công bố kết quả khoa học.}
    Hoàn thiện số liệu thực nghiệm sim-to-real và chuẩn bị bài báo gửi
    tạp chí / hội nghị kỹ thuật trong nước phù hợp (ví dụ: VCCA, REV).
\end{enumerate}
